%%% FIELD proof-bounded-overlap-second-order
Fix \(P\in\mathcal P(d,0,\alpha,\beta,\gamma,L)\).  By the definition of \(r_{\kappa}\), \(r_0(x)=1\) for every \(x\in\mathcal X\).  Therefore \(\mathrm{ass:local\mbox{-}propensity\mbox{-}shape}\) and \(\mathrm{ass:multiplier\mbox{-}envelope}\) give
\[
 e_P(x)=\frac{r_0(x)g_P(x)}4=\frac{g_P(x)}4,
 \qquad
 \frac14\leq e_P(x)\leq\frac12
 \tag{35}
\]
for every \(x\in\mathcal X\).  Thus this is the original exact-uniform model on its bounded-overlap parameter slice, not a replacement class.

For the lower bound, every data-only Borel estimator \(T\) induces the admissible supplied-propensity estimator \(T^e(o_{1:n},e)=T(o_{1:n})\), which ignores its second argument and has the same loss under every law.  The definitions \(\mathrm{def:unknown\mbox{-}risk}\) and \(\mathrm{def:oracle\mbox{-}risk}\) therefore imply
\[
 R_n^e(d,0,\beta,\gamma,L)
 \leq R_n(d,0,\alpha,\beta,\gamma,L).
 \tag{36}
\]
The supplied-propensity lower bound in \(\mathrm{thm:oracle\mbox{-}minimax}\), specialized to \(\kappa=0\), and (36) give
\[
 c_e n^{-\gamma/(2\gamma+d)}
 \leq R_n(d,0,\alpha,\beta,\gamma,L)
 \qquad(n\geq2).
 \tag{37}
\]
The two laws used there are genuine full-support uniform-design Bernoulli laws in the exact class.

For the upper bound, invoke the already established \(\mathrm{lem:bounded\mbox{-}overlap\mbox{-}second\mbox{-}order\mbox{-}risk}\).  Put \(S=\alpha+\beta\) and \(\ell=hk^{-1/d}\).  Since \(h\in(0,1]\), \(k\geq1\), and \(\alpha>0\), one has \(0<\ell\leq1\) and hence \(h^\gamma\ell^\alpha\leq h^\gamma\).  The lemma's explicitly defined estimator thus satisfies
\[
 \sup_{P\in\mathcal P(d,0,\alpha,\beta,\gamma,L)}
 \mathbb E_P\left|\widehat\tau^{(2)}_{n,h,k}(0)-\tau_P(0)\right|
 \leq C\left\{h^\gamma+\ell^S+(nh^d)^{-1/2}
 +\frac{\sqrt{k}}{nh^d}\right\}.
 \tag{38}
\]

First suppose \(S<S_\star\).  Put \(\delta=D_\star^{-1}\) and choose, up to the bounded-ratio admissible-rank rounding built into the lemma,
\[
 h_n=n^{-1/(\gamma D_\star)}=n^{-\delta/\gamma},
 \qquad
 k_n\asymp n^{(d/S-d/\gamma)/D_\star}
 =n^{\delta d(1/S-1/\gamma)}.
 \tag{39}
\]
Because \(S<S_\star<\gamma\), the exponent of \(k_n\) is positive.  With \(\ell_n=h_nk_n^{-1/d}\), equation (39) yields
\[
 \ell_n\asymp n^{-\delta/S},
 \qquad
 h_n^\gamma\asymp\ell_n^S\asymp n^{-\delta}.
 \tag{40}
\]
Moreover, the identity \(1=\delta\{1+d/(2\gamma)+d/(2S)\}\) gives
\[
 \frac{\sqrt{k_n}}{nh_n^d}
 \asymp n^{-1+\delta d/(2S)+\delta d/(2\gamma)}
 =n^{-\delta}.
 \tag{41}
\]
The exponent of the nondegenerate term \((nh_n^d)^{-1/2}\) exceeds \(\delta\) by
\[
 \left(\frac12-\frac{\delta d}{2\gamma}\right)-\delta
 =\frac{d(1/S-1/\gamma)-2}{4D_\star}\geq0,
 \tag{42}
\]
where the last inequality is equivalent to \(S\leq d\gamma/(2\gamma+d)=S_\star\).  Substitution of (40)--(42) into (38) proves
\[
 R_n(d,0,\alpha,\beta,\gamma,L)\leq Cn^{-1/D_\star}
 \qquad(S<S_\star).
 \tag{43}
\]

Now suppose \(S\geq S_\star\).  Choose
\[
 h_n=n^{-1/(2\gamma+d)},
 \qquad
 k_n\asymp nh_n^d,
 \tag{44}
\]
again using the lemma's deterministic rank rounding.  Then
\[
 h_n^\gamma=(nh_n^d)^{-1/2}asymp\frac{\sqrt{k_n}}{nh_n^d}
 =n^{-\gamma/(2\gamma+d)},
 \qquad
 \ell_n\asymp n^{-1/d}.
 \tag{45}
\]
The definition of \(S_\star\) gives the exact equivalence
\[
 S\geq S_\star
 \quad\Longleftrightarrow\quad
 \frac Sd\geq\frac\gamma{2\gamma+d},
\]
so (45) implies
\[
 \ell_n^S\lesssim n^{-\gamma/(2\gamma+d)}.
 \tag{46}
\]
Equations (38) and (44)--(46) prove the supplied-propensity-rate upper bound for \(S\geq S_\star\).  The estimator from the lemma is total for every \(n\geq2\); deterministic rounding and the finitely many small-sample cases are absorbed by enlarging \(C\), since clipping bounds the absolute loss by \(3/2\).

Finally, \(D_\star^{-1}\leq\gamma/(2\gamma+d)\) holds exactly when \(S\leq S_\star\).  Thus (43) and the high-smoothness upper bound combine to give
\[
 R_n(d,0,\alpha,\beta,\gamma,L)
 \leq C\max\left\{n^{-\gamma/(2\gamma+d)},n^{-1/D_\star}\right\}.
 \tag{47}
\]
Together with (37), this proves the asserted bracket and rate equality when \(S\geq S_\star\).  At \(S=S_\star\),
\[
 D_\star=2+\frac d\gamma,
 \qquad
 \frac1{D_\star}=\frac\gamma{2\gamma+d}.
 \tag{48}
\]
Equations (45), (46), and (48) show that finitely many pure-power bias and variance terms agree; no dyadic or harmonic summation occurs.  Hence no logarithmic multiplier is introduced at equality.  For \(S<S_\star\), (37) and (43) give only the stated bracket, and no nuisance-dominated converse has been inferred.
%%% FIELD proof-oracle-minimax
Write \(m_\gamma=\lceil\gamma\rceil-1\), let \(v_\gamma(z)\) be the finite vector of all monomials \(z^\nu\) with \(|\nu|\leq m_\gamma\), and put \(D=[1/2,1]^d\).  Its Gram matrix
\[
 H_\gamma=\int_Dv_\gamma(z)v_\gamma(z)^{\mathsf T}\,dz
\]
is positive definite: if \(c^{\mathsf T}H_\gamma c=0\), then the continuous polynomial \(c^{\mathsf T}v_\gamma\) vanishes almost everywhere on \(D\), hence everywhere on the interior of \(D\), and therefore all its coefficients vanish.  Define
\[
 q_\gamma(z)=v_\gamma(0)^{\mathsf T}H_\gamma^{-1}v_\gamma(z).
\]
It follows directly that, for every polynomial \(p\) of total degree at most \(m_\gamma\),
\[
 \int_Dq_\gamma(z)p(z)\,dz=p(0).
 \tag{1}
\]

The condition \(\mathrm{ass:contrast\mbox{-}holder}\), applied on each segment from \(0\) to \(x\), gives a polynomial \(p_{P,h}\) of degree at most \(m_\gamma\) such that
\[
 \sup_{z\in D}\left|\tau_P(hz)-p_{P,h}(z)\right|\leq C_1h^\gamma,
 \qquad p_{P,h}(0)=\tau_P(0).
 \tag{2}
\]
For completeness, when \(m_\gamma=0\), this is precisely the defining H\"older inequality.  When \(m_\gamma\geq1\), apply the fundamental theorem of calculus successively to \(t\mapsto\tau_P(thz)\) through order \(m_\gamma\).  After subtracting the Taylor polynomial, the defining \((\gamma-m_\gamma)\)-H\"older bound makes the increment of every order-\(m_\gamma\) derivative at most \(L\lVert hz\rVert_\infty^{\gamma-m_\gamma}\), while the remaining \(m_\gamma\) integrations contribute \(\lVert hz\rVert_\infty^{m_\gamma}\).  Summing over the finitely many multi-indices proves (2), including integer \(\gamma\), for which the order-\((\gamma-1)\) derivatives are Lipschitz.

For \(h\in(0,1]\), let \(D_h=hD\).  The supplied propensity is never inverted outside \(D_h\).  On \(D_h\), the definition of \(r_\kappa\) and the propensity-shape and multiplier-envelope properties incorporated in \(\mathrm{def:exact\mbox{-}model}\) give
\[
 \frac{h^\kappa}{2^\kappa4}\leq e_P(x)\leq\frac{h^\kappa}{2},
 \qquad 1-e_P(x)\geq\frac12,
 \tag{3}
\]
where the first lower bound is also valid at \(\kappa=0\).  Define
\[
 Z_P(O)=\frac{AY}{e_P(X)}-\frac{(1-A)Y}{1-e_P(X)}
\]
on \(\{X\in D_h\}\), and set the summand below to zero off that set.  By \(\mathrm{ass:bernoulli\mbox{-}response}\),
\[
 \mathbb E_P[Z_P(O)\mid X=x]=\mu_{1,P}(x)-\mu_{0,P}(x)=\tau_P(x),
 \tag{4}
\]
and, because \(Y\) is binary, using (3) and the disjointness of the two treatment indicators gives
\[
 \mathbb E_P[Z_P(O)^2\mid X=x]
 =\frac{\mu_{1,P}(x)}{e_P(x)}+\frac{\mu_{0,P}(x)}{1-e_P(x)}
 \leq C_2h^{-\kappa}.
 \tag{5}
\]
The total oracle estimator is
\[
 \widehat\tau^{e}_{h}
 =\Pi_{[-3/4,3/4]}\!\left[
 \frac1n\sum_{i=1}^nh^{-d}q_\gamma(X_i/h)
 \mathbf 1\{X_i\in D_h\}Z_P(O_i)
 \right],
 \tag{6}
\]
where \(\Pi_I(u)=\min\{\max\{u,\inf I\},\sup I\}\).  Formula (6) is Borel and defined on every sample; (3) shows that it never divides by zero.  By the uniform-design property incorporated in \(\mathrm{def:exact\mbox{-}model}\), (1), (2), and (4), its unclipped expectation has bias at most
\[
 \left|\int_Dq_\gamma(z)\tau_P(hz)\,dz-\tau_P(0)\right|
 \leq \lVert q_\gamma\rVert_{L^1(D)}C_1h^\gamma.
 \tag{7}
\]
By \(\mathrm{ass:iid\mbox{-}sampling}\), (5), and the change of variables \(x=hz\),
\[
 \operatorname{Var}_P(\widehat\tau^{e}_{h,\mathrm{unclipped}})
 \leq\frac1n\int_{D_h}h^{-2d}q_\gamma(x/h)^2C_2h^{-\kappa}\,dx
 \leq\frac{C_3}{nh^{d+\kappa}}.
 \tag{8}
\]
Clipping cannot increase distance to \(\tau_P(0)\in[-3/4,3/4]\).  Therefore (7), (8), and \(\mathbb E|W-\mathbb EW|\leq\{\operatorname{Var}(W)\}^{1/2}\) yield
\[
 \sup_{P\in\mathcal P(d,\kappa,\alpha,\beta,\gamma,L)}
 \mathbb E_P|\widehat\tau_h^e-\tau_P(0)|
 \leq C_4\left\{h^\gamma+(nh^{d+\kappa})^{-1/2}\right\}.
 \tag{9}
\]
Taking \(h=c_0n^{-1/(2\gamma+d+\kappa)}\), with a fixed \(c_0\in(0,1]\), proves the asserted oracle upper bound.

For the converse, choose a fixed \(C^\infty\) function \(\psi\) on the cube with \(\psi(0)=1\), \(|\psi|\leq1\), and support contained in \([0,1/2]^d\).  Fix \(g_P\equiv1\), \(e_P=r_\kappa/4\), and \(\mu_{0,P}\equiv1/2\).  For \(\varepsilon\in\{-1,1\}\), set
\[
 \tau_\varepsilon(x)=\varepsilon a h^\gamma\psi(x/h),
 \qquad \mu_{1,\varepsilon}(x)=\frac12+\tau_\varepsilon(x),
 \tag{10}
\]
where \(a>0\) is fixed sufficiently small as a function only of \((d,\gamma,L)\).  Scaling each derivative in the definition of \(\mathcal H_d^\gamma(L)\) shows that the H\"older seminorm of the perturbation in (10) is at most a constant times \(a\); its lower-order derivative suprema are at most the same constant because \(h\leq1\).  Thus (10), the constant baseline, and the constant multiplier obey all H\"older restrictions.  Decreasing \(a\) if necessary gives \(1/8\leq\mu_{0,P},\mu_{1,P}\leq7/8\).

These observed laws have genuine causal completions.  Conditional on \(X=x\), draw \(Y(0)\) and \(Y(1)\) independently as Bernoulli variables with means \(\mu_{0,P}(x)\) and \(\mu_{1,\varepsilon}(x)\), draw \(A\) independently of them with mean \(e_P(x)\), and set \(Y=AY(1)+(1-A)Y(0)\).  Consistency and conditional exchangeability then hold by construction.  Hence \(\mathrm{def:exact\mbox{-}model}\) shows that both laws, denoted \(P_+\) and \(P_-\), belong to the exact class and supply the oracle with the same function \(e_P\).

For Bernoulli probabilities \(u,v\in[1/8,7/8]\), direct summation gives
\[
 \chi^2\{\operatorname{Bernoulli}(u),\operatorname{Bernoulli}(v)\}
 =\frac{(u-v)^2}{v(1-v)}\leq C_5(u-v)^2.
 \tag{11}
\]
The treatment and control mechanisms are otherwise identical, so (10), (11), and the uniform-design property in \(\mathrm{def:exact\mbox{-}model}\) imply
\[
 \chi^2(P_+,P_-)
 \leq C_6a^2h^{2\gamma}\int_{[0,h/2]^d}r(x)^\kappa\,dx
 \leq C_7a^2h^{2\gamma+d+\kappa}.
 \tag{12}
\]
The last power follows from \(x=hz\); the remaining integral is finite and positive for every \(\kappa\geq0\).  Independence gives the exact identity
\[
 1+\chi^2(P_+^{\otimes n},P_-^{\otimes n})
 =\{1+\chi^2(P_+,P_-)\}^n.
 \tag{13}
\]
Choose the fixed constant \(c_0\) in \(h=c_0n^{-1/(2\gamma+d+\kappa)}\) sufficiently small that (12)--(13) make the product chi-square divergence at most \(1/4\).  If \(Q\ll R\), Cauchy--Schwarz applied to \(\int|dQ/dR-1|\,dR\) yields
\[
 \lVert Q-R\rVert_{\mathrm{TV}}\leq\frac12\sqrt{\chi^2(Q,R)};
\]
hence the product total variation is at most \(1/4\).

Finally, put \(\Delta=|\tau_+(0)-\tau_-(0)|=2ah^\gamma\).  For any oracle estimator \(T^e\), the two events \(E_\varepsilon=\{|T^e-\tau_\varepsilon(0)|\geq\Delta/2\}\) cover the sample space pointwise.  Integrating their indicators against the smaller of the two product densities gives
\[
 P_+^{\otimes n}(E_+)+P_-^{\otimes n}(E_-)
 \geq1-\lVert P_+^{\otimes n}-P_-^{\otimes n}\rVert_{\mathrm{TV}}
 \geq\frac34.
\]
Consequently,
\[
 \max_{\varepsilon\in\{-1,1\}}
 \mathbb E_{P_\varepsilon}|T^e-\tau_\varepsilon(0)|
 \geq\frac{\Delta}{4}\cdot\frac34
 \geq c_1h^\gamma.
\]
Taking the infimum over oracle estimators proves the lower bound.  Enlarging constants covers every \(n\geq2\).  Every constant is independent of \(n\); neither construction uses \(\alpha\), and the constant baseline belongs to every prescribed \(\beta\)-ball.
